% =====================================================================
\phantomsection
\addcontentsline{toc}{part}{Compiled Languages for Speed and Systems}
\BookChapter[Fortran]{Fortran: Formulas at Full Speed}
\label{ch:fortran}

\begin{researchquestion}
Why is a language from 1957 still running weather forecasts, crash simulations and materials calculations?
\end{researchquestion}
\OrientationNote
% =====================================================================

\begin{factsheet}{Fortran}
First appeared & 1957 (FORTRAN I); modern standards Fortran 90, 95, 2003, 2008, 2018, 2023\\
Designed by & John Backus and team, IBM\\
Paradigms & procedural, array programming; object-oriented since Fortran 2003; parallel (coarrays)\\
Typing & static, strong (with \code{implicit none})\\
Execution & compiled to native code\\
File extension & \file{.f90} (free form); \file{.f}, \file{.for} (old fixed form)\\
Compilers & gfortran (free, GNU), Intel ifx (free), LLVM Flang (free), NVIDIA nvfortran, Cray\\
Used for & weather and climate models, computational fluid dynamics, structural analysis,
  quantum chemistry and materials simulation, nuclear engineering, numerical libraries\\
Famous software & LAPACK, BLAS, VASP, Quantum ESPRESSO, WRF (weather), NASTRAN (origins),
  many CFD and reactor codes, parts of SciPy\\
\end{factsheet}

Fortran is the oldest programming language still in wide use, and it is still one of the
fastest. If you work in computational mechanics, CFD, climate science or materials
simulation, you will almost certainly meet it. Modern Fortran looks very different from
the FORTRAN 77 code of the 1980s that gave the language a reputation for being
old-fashioned: it has free-form source, modules, dynamic arrays, whole-array operations,
objects and built-in parallelism.

\LearningOutcomes{Compile and run a modern Fortran program; declare variables with an explicit precision; use whole-array operations, functions, subroutines and modules; read formatted output statements; recognise old fixed-form FORTRAN 77 code; and explain how Fortran and Python are used together.}

\section{History and philosophy}

\begin{history}
In 1953 John Backus proposed to his managers at IBM that a program should translate
mathematical formulas into machine code automatically. Sceptics argued that such
translated code would be far slower than hand-written machine code and that no serious
programmer would use it. Backus's team spent three years building an \emph{optimising}
compiler for the IBM 704 and released FORTRAN in 1957. The compiled code was nearly as fast
as hand-written code, and within a few years most scientific programming in the world was
done in FORTRAN. Backus later received the Turing Award, the ``Nobel prize'' of computing.
\end{history}

Fortran's design priorities have stayed the same for almost 70 years:
\begin{enumerate}
\item \textbf{Numerical performance.} The language is designed so that compilers can
  optimise array loops aggressively. For example, Fortran assumes that two array arguments
  of a subroutine do not overlap in memory (no \emph{aliasing}), which lets the compiler
  vectorise loops that a C compiler must treat cautiously.
\item \textbf{Arrays as first-class objects.} Multidimensional arrays, array slices and
  whole-array arithmetic are built into the language, not added by a library.
\item \textbf{Backward compatibility.} A FORTRAN 77 program from 1985 usually still compiles
  today. This is why so much valuable, carefully validated engineering code is still in
  use: rewriting it would be expensive and risky.
\end{enumerate}

The upper-case spelling FORTRAN is used for the versions up to FORTRAN 77; since Fortran 90
the official name is written ``Fortran''.

\section{Getting started}

Install the free \code{gfortran} compiler (part of GCC): on Linux with the package
manager, on macOS with Homebrew (\code{brew install gcc}), and on Windows through MSYS2,
WSL, or the Intel oneAPI toolkit. The \emph{Fortran Package Manager} (\code{fpm}) and the
website \url{https://fortran-lang.org} are good starting points. Online compilers such as
Compiler Explorer (godbolt.org) let you try Fortran in a browser.
\codefile{[08]Fortran}{fortran/hello.f90}{hello world}
\begin{lstlisting}[style=shell]
$ gfortran -o hello hello.f90
$ ./hello
\end{lstlisting}
\outfile{fortran/hello.out}

Every Fortran program starts with \code{program name} and ends with \code{end program
name}. The statement \code{implicit none} should be the first line of every program unit.
It switches off a historical rule (\term{implicit typing}) by which variables whose names
began with the letters I to N were automatically integers and all others were reals. Under
that rule a misspelt variable silently became a new variable --- a famous source of bugs.
\code{print *,} writes values to the screen in a default (``list-directed'') format.

Fortran is \emph{not case sensitive}: \code{PRINT}, \code{Print} and \code{print} are the
same, and so are \code{Velocity} and \code{velocity}. Comments start with \code{!}. A
statement that is too long for one line is continued with an ampersand \code{\&} at the end
of the line.

\section{Variables and types}

All variables must be declared at the beginning of a program unit, before any executable
statement:
\begin{lstlisting}[language={[08]Fortran}]
integer :: n_steps, i
real    :: temperature              ! default real: usually single precision!
real(8) :: pressure                 ! common but non-portable way to get double
logical :: converged
character(len=20) :: material
complex :: z
\end{lstlisting}
The intrinsic types are \code{integer}, \code{real}, \code{complex}, \code{logical}
(\code{.true.} or \code{.false.}) and \code{character}. Constants are declared with the
\code{parameter} attribute and cannot be changed:
\begin{lstlisting}[language={[08]Fortran}]
real(dp), parameter :: g = 9.81_dp, pi = 4.0_dp * atan(1.0_dp)
\end{lstlisting}

\subsection*{Choose the precision with a kind parameter}
\label{sec:fortran-kinds}

The default \code{real} in Fortran is \emph{single precision} (about 7 digits). For
engineering work you almost always want double precision. The portable way to ask for it
is to define a \term{kind} parameter once and use it everywhere:
\begin{lstlisting}[language={[08]Fortran}]
integer, parameter :: dp = selected_real_kind(15, 307)   ! >= 15 digits, 10^307
! or, equivalently on all common compilers:
use, intrinsic :: iso_fortran_env, only: dp => real64
real(dp) :: x
x = 0.1_dp        ! the suffix _dp makes the literal double precision
\end{lstlisting}

\begin{pitfall}[Common pitfall: single-precision literals]
In Fortran, \code{0.1} is a \emph{single-precision} constant even if you assign it to a
double-precision variable. \code{x = 0.1} stores $0.100000001490116\ldots$, accurate to only
7 digits. Always write \code{0.1\_dp} (or \code{0.1d0} in older code). This is one of the
most common precision bugs in Fortran programs.
\end{pitfall}

\subsection*{Operators look like those of Python, with a few traps}

Arithmetic uses \code{+ - * /} and \code{**} for powers, as in \py. Integer division
truncates: \code{7/2} is \code{3}. Comparisons can be written in modern form
(\code{== /= < <= > >=}) or in the old form (\code{.eq. .ne. .lt. .le. .gt. .ge.}); note
that ``not equal'' is \code{/=}, not \code{!=}. Logical operators are \code{.and.},
\code{.or.} and \code{.not.}. Strings are joined with \code{//}.

Fortran has a rich set of \term{intrinsic} (built-in) functions: \code{sqrt}, \code{exp},
\code{log}, \code{log10}, \code{sin}, \code{cos}, \code{tan}, \code{atan2}, \code{abs},
\code{mod}, \code{min}, \code{max}, \code{nint} (nearest integer), \code{real} and
\code{int} (conversions), and many array functions described below.

\section{Control flow}

\begin{lstlisting}[language={[08]Fortran}]
! if / else if / else
if (Re < 2300.0_dp) then
    regime = "laminar"
else if (Re < 4000.0_dp) then
    regime = "transitional"
else
    regime = "turbulent"
end if

! counted loop: i = 1, 2, ..., 10 (the end value IS included)
do i = 1, 10
    total = total + 1.0_dp / i**2
end do

! loop with a step
do i = 10, 1, -2        ! 10, 8, 6, 4, 2
    print *, i
end do

! loop until a condition is met
do
    x = x - f(x) / df(x)
    if (abs(f(x)) < 1.0e-12_dp) exit      ! exit = break
end do

! select case: a clean multi-way branch
select case (n_dims)
case (1)
    print *, "line"
case (2)
    print *, "surface"
case default
    print *, "volume"
end select
\end{lstlisting}
\code{exit} leaves a loop (like \code{break} in \py) and \code{cycle} jumps to the next
iteration (like \code{continue}). Old code is full of \code{GOTO} statements and numbered
labels; avoid them in new code.

\section{Arrays}

Arrays are where Fortran shines. By default, indices start at 1, but you may choose any
range:
\begin{lstlisting}[language={[08]Fortran}]
real(dp) :: x(100)            ! x(1) ... x(100)
real(dp) :: A(3, 3)           ! 3 x 3 matrix
real(dp) :: v(0:40)           ! indices 0 ... 40
real(dp), allocatable :: u(:) ! size decided at run time
allocate(u(n))                ! ... and released automatically at the end of scope
\end{lstlisting}
Slices use the notation \code{start:end:stride}, and \emph{the end is included}:
\code{x(2:6:2)} means $x_2, x_4, x_6$. Whole-array expressions work element by element,
like NumPy or MATLAB, and array constructors use square brackets:
\codefile{[08]Fortran}{fortran/arrays.f90}{whole-array operations}
\outfile{fortran/arrays.out}
Note the \code{reshape} call: because Fortran stores arrays in column-major order
(Section~\ref{sec:rowcolumn}), the list of values fills the matrix \emph{column by
column}. Here the matrix is symmetric, so it does not matter, but it often does when you
pass arrays between Fortran and C or \py. Useful array intrinsics include \code{sum},
\code{product}, \code{maxval}, \code{minval}, \code{maxloc}, \code{count}, \code{any},
\code{all}, \code{size}, \code{shape}, \code{matmul}, \code{dot\_product},
\code{transpose} and \code{reshape}. The \code{where} construct applies an assignment only
where a mask is true:
\begin{lstlisting}[language={[08]Fortran}]
where (T > 100.0_dp)
    phase = 1        ! vapour
elsewhere
    phase = 0        ! liquid
end where
\end{lstlisting}

\section{Procedures: functions and subroutines}

Fortran has two kinds of procedures. A \term{function} returns one value and is used in
expressions. A \term{subroutine} returns nothing directly but can modify its arguments;
it is invoked with \code{call}. Every argument should declare its \term{intent}:
\code{in} (read only), \code{out} (written) or \code{inout} (both). The compiler then
catches attempts to modify an input.
\begin{lstlisting}[language={[08]Fortran}]
subroutine stats(x, mean, std)
    real(dp), intent(in)  :: x(:)          ! assumed-shape array
    real(dp), intent(out) :: mean, std
    integer :: n
    n = size(x)
    mean = sum(x) / n
    std = sqrt(sum((x - mean)**2) / (n - 1))
end subroutine stats

! ... in the main program:
call stats(yield_strength, m, s)
\end{lstlisting}
Unlike C or \py, Fortran passes arguments \emph{by reference}: the subroutine works on the
caller's variable, not on a copy. This is efficient for large arrays and is why
\code{intent} matters.

Procedures marked \code{pure} have no side effects, and \code{elemental} procedures are
written for scalars but automatically apply element by element to arrays, like a NumPy
\emph{ufunc}. Procedures can be placed after a \code{contains} statement inside the main
program (as in the drag example below) or, better, in a module.

\section{Modules}

A \term{module} groups constants, types and procedures that belong together, and controls
what is visible from outside with \code{public} and \code{private}. Other program units
access it with \code{use}. Modules also let the compiler check every call against the
procedure's declared arguments, which catches many errors at compile time. The following
example computes the deflection of a cantilever beam with an end load $P$,
\[
w(x) = \frac{P x^2 (3L - x)}{6EI},
\]
using an \code{elemental} function so that it works on a whole array of positions at once.
\codefile{[08]Fortran}{fortran/beam_module.f90}{a module and an elemental function}
\outfile{fortran/beam_module.out}
The expression \code{[(L * k / 4.0\_dp, k = 0, 4)]} is an \term{implied-do} array
constructor, Fortran's equivalent of a list comprehension. In a larger project each module
lives in its own file, and the build system compiles modules before the files that use
them.

\section{Input and output}

Fortran's I/O uses \emph{unit numbers} and \emph{format descriptors}. \code{write(*, fmt)}
writes to the screen; \code{read(*, *)} reads from the keyboard. Format descriptors play
the same role as the format specifications in \py's f-strings:
\begin{center}
\small
\begin{tabular}{@{}lll@{}}
\toprule
\textbf{Descriptor} & \textbf{Meaning} & \textbf{\py equivalent}\\
\midrule
\code{i5} & integer, width 5 & \code{\{n:5d\}}\\
\code{f10.4} & fixed point, width 10, 4 decimals & \code{\{x:10.4f\}}\\
\code{es12.4} & scientific notation & \code{\{x:12.4e\}}\\
\code{a} & character string & \code{\{s\}}\\
\code{3f8.2} & three reals & three fields\\
\code{f0.4} & minimal width, 4 decimals & \code{\{x:.4f\}}\\
\bottomrule
\end{tabular}
\end{center}
Files are opened with an explicit unit:
\begin{lstlisting}[language={[08]Fortran}]
integer :: u, ios, n
real(dp) :: value, total
open(newunit=u, file="yield.txt", status="old", action="read")
total = 0.0_dp;  n = 0
do
    read(u, *, iostat=ios) value
    if (ios /= 0) exit            ! end of file (or a read error)
    total = total + value
    n = n + 1
end do
close(u)
print '(a, i0, a, f8.2)', "Read ", n, " values, mean = ", total / n
\end{lstlisting}
For large binary data sets, Fortran codes often use libraries such as HDF5 or NetCDF,
which can also be read from \py.

\section{Case study: the falling skydiver}

\codefile{[08]Fortran}{fortran/drag.f90}{Euler solution of the drag equation}
\begin{lstlisting}[style=shell]
$ gfortran -O2 -Wall -o drag drag.f90
$ ./drag
\end{lstlisting}
\outfile{fortran/drag.out}
The numbers agree exactly with the \py output in Chapter~\ref{ch:python},
as they must: both programs perform the same IEEE double-precision arithmetic in the same
order. Points to notice:
\begin{itemize}
\item the kind parameter \code{dp} and the \code{\_dp} suffix on every literal;
\item the array \code{v(0:n\_steps)} indexed from 0, so that \code{v(n)} is the speed at
  time $n\,\Delta t$;
\item the functions placed after \code{contains} can use the constants \code{m}, \code{c}
  and \code{g} of the main program (\term{host association});
\item \code{pure} functions with \code{intent(in)} arguments.
\end{itemize}

\section{Object-oriented and parallel Fortran}

Fortran 2003 added \term{derived types} with type-bound procedures, inheritance
(\code{extends}) and polymorphism:
\begin{lstlisting}[language={[08]Fortran}]
type :: pipe_t
    real(dp) :: length, diameter
contains
    procedure :: area => pipe_area
end type pipe_t
...
type(pipe_t) :: p
p = pipe_t(12.0_dp, 0.05_dp)
print *, p%area()            ! % selects a component, like . in Python
\end{lstlisting}
For parallel computing, Fortran codes typically use \textbf{OpenMP} (directives that
split loops over the cores of one computer) and \textbf{MPI} (message passing between many
computers in a cluster). Fortran 2008 also introduced \term{coarrays}, a parallel model
built into the language itself, and \code{do concurrent}, which tells the compiler that
loop iterations are independent:
\begin{lstlisting}[language={[08]Fortran}]
!$omp parallel do
do i = 1, n
    y(i) = a * x(i) + y(i)
end do
!$omp end parallel do
\end{lstlisting}

\section{Fortran and Python together}

Many engineers use Fortran for the fast core of a calculation and \py for everything
around it. Tools for connecting the two include:
\begin{itemize}
\item \textbf{f2py}, part of NumPy, which wraps Fortran subroutines so that they can be
  imported as a \py module;
\item Fortran's standard \code{iso\_c\_binding} module, which makes Fortran procedures
  callable from C, and therefore from \py (via \code{ctypes} or \code{cffi}), Julia and
  others;
\item running a compiled Fortran program from \py with \code{subprocess} and reading its
  output files --- the most common approach in practice.
\end{itemize}

\begin{engnote}[Engineering note: reading legacy code]
You may be asked to modify a FORTRAN 77 program. Recognise its features: code starting in
column 7 (columns 1--5 for labels, column 6 for continuation), comments marked by \code{C}
or \code{*} in column 1, \code{COMMON} blocks for shared global data, \code{DO 10 I=1,N}
loops ending at a labelled \code{10 CONTINUE}, \code{GOTO} jumps and implicit typing.
Change as little as possible, add \code{IMPLICIT NONE} if you can, and above all build a
set of test cases \emph{before} you change anything, so that you can show the results did
not change.
\end{engnote}

\section{Strengths and weaknesses}

\begin{center}
\small
\begin{tabularx}{\linewidth}{@{}>{\raggedright\arraybackslash}X>{\raggedright\arraybackslash}X@{}}
\toprule
\textbf{Strengths} & \textbf{Weaknesses}\\
\midrule
Excellent performance for numerical array code & Small community compared with \py or C++\\
Arrays, slices and complex numbers built in & Weak string handling, few general-purpose libraries\\
Very stable standard; decades-old code still compiles & Large amounts of old-style legacy code\\
Mature parallel computing (OpenMP, MPI, coarrays) & Fewer modern tools (package manager is young)\\
Simple language for formulas; hard to write very slow code & Not suited to user interfaces, web or embedded work\\
\bottomrule
\end{tabularx}
\end{center}

\section{Further reading}
\begin{itemize}
\item \url{https://fortran-lang.org} --- tutorials, the Fortran Package Manager and the
  Fortran standard library project.
\item M.~Metcalf, J.~Reid and M.~Cohen, \emph{Modern Fortran Explained: Incorporating
  Fortran 2023}, Oxford University Press, 2023.
\item M.~Curcic, \emph{Modern Fortran: Building Efficient Parallel Applications},
  Manning, 2020.
\item S.~J.~Chapman, \emph{Fortran for Scientists and Engineers}, 4th ed., McGraw-Hill,
  2017.
\end{itemize}

\begin{exercises}
\item Compile and run \file{drag.f90}. Then remove the \code{\_dp} suffix from the
  constant \code{c = 0.25\_dp} and from \code{g = 9.81\_dp}. Does the output change? Explain
  using the pitfall box in Section~\ref{sec:fortran-kinds}.
\item Write a Fortran program that prints a table of temperatures from $-20$ to
  $100\,^\circ$C in steps of 20 together with their values in Fahrenheit and kelvin, using
  formatted output.
\item Write a subroutine \code{stats(x, mean, std)} in a module and use it to compute the
  mean and standard deviation of the values in \file{yield.txt}. Compare with the \py
  result in Chapter~\ref{ch:python}.
\item Use \code{matmul} to verify that the vector $\mathbf{x} = (1, -2, 3)$ solves the
  system in \file{arrays.f90}. Then look up the LAPACK routine \code{dgesv} and try to call
  it to solve the system (compile with \code{-llapack}).
\item Translate this FORTRAN 77 fragment into modern Fortran:
\begin{lstlisting}[language={[77]Fortran}]
      DIMENSION X(100)
      S = 0.0
      DO 10 I = 1, N
         S = S + X(I)**2
   10 CONTINUE
      RMS = SQRT(S/N)
\end{lstlisting}
  Which variables are integers and which are reals under implicit typing? Can you replace
  the loop by a single array expression?
\item Write a loop that sums the elements of a $2000\times2000$ matrix, once with the first
  index in the inner loop and once with the second. Time both with \code{cpu\_time}. Which is
  faster, and why?
\end{exercises}
